% Source: source_md/intelligence_part2_combined_ja.md
\section{ゲティア問題 — $C \neq$ 真理の証拠、$L_E$ の操作的安定性}


\S{}6〜\S{}9 で、停止が真理によってではなく構造的安定性と改訂可能性によって維持されることを見た。本章は、この事実が認識論において最も鋭く現れる場所——ゲティア問題——を分析する。ゲティア問題は、閉包度 $C$ が真理測度ではなく安定性測度であることの構造的証拠である。

本章は Part I \S{}6(ゲティア問題の構造的診断)を継承し、三層動力学の語彙でこれを再展開する。

\bigskip


\subsection{ゲティア問題は構造的事実の暴露である}


ゲティア問題は、伝統的に「知識とは正当化された真なる信念である」という定義への反例として知られる。正当化され、真であり、信じられているにもかかわらず、知識とは見なしがたい事例が構成できる——これがゲティア問題である。

本論文の立場では、ゲティア問題は周辺的な反例の寄せ集めではない。それは構造的事実の暴露である。すなわち、正当化され真である信念が、なお知識として安定しないことがあるという事実は、知識が正当化と真理の合成によっては定義できないことを示している。

三層動力学の語彙で言えば、ゲティア問題が暴露するのは次のことである。推論層 $L_F$ における正当化(\S{}3 の正当化測度 $J$)と、外部実在に対応する真理は、創発停止層 $L_E$ における停止とは別の事柄である。$L_E$ が供給する停止は、正当化と真理の合致によってではなく、構造的安定性によって成立する。ゲティア事例は、正当化と真理が合致してもなお $L_E$ の安定が独立の事柄であることを示す。

\bigskip


\subsection{ゲティア問題は原理的に解決不可能である}


ゲティア問題に対しては、長年にわたり「第四の条件」を追加して知識を再定義する試みがなされてきた。本論文の立場では、これらの試みは原理的に成功しない。

理由は \S{}3 と \S{}5 から導かれる。知識を「正当化+真理+α」として定義しようとする試みは、α として何らかの条件を $L_F$ の内部に求める。しかし \S{}3 で示した通り、$L_F$ は内部に停止条件を持たない。知識を成立させる「最終的な正当化条件」を $L_F$ 内部に定義しようとすると、それは内部停止述語の定義に帰着し、無限後退する(系 3.2)。

したがって、ゲティア問題を $L_F$ 内部の条件追加によって解決することは原理的にできない。知識の成立は $L_F$ 内部では閉じない。それは $L_E$——外部依存的な創発停止層——においてのみ成立する。ゲティア問題が解決不可能なのは、知識の成立条件が $L_F$ 内部にないからである。

これは欠陥ではなく、\S{}3 の非閉包定理の認識論的帰結である。ゲティア問題の解決不可能性は、知識が個体内の正当化では閉じないことの、認識論における現れである。

\bigskip


\subsection{社会はゲティア問題を解決しない}


ゲティア問題に対する社会的停止構造($\Phi_{\text{net}}$)の態度は、注目に値する。社会はゲティア問題を解決しない。そして、解決しようとしない。

\S{}9 で見た通り、社会的ブロックチェーンは正統性によって停止を供給する。正統性は計算不可能であり(\S{}9.4)、完全な正当化を要求しない。社会は、ある判断が完全に正当化され真理に対応することを確認してから受容するのではない。社会は、暫定的に受容可能な判断をブロックとして採用し、改訂可能に保つ。

もし社会が完全な正当化を停止の条件として要求すれば、何が起こるか。完全な正当化は \S{}3 により到達不可能である。したがって、完全な正当化を要求する社会的停止構造は、いかなる判断も受容できず、即座に停止する。社会的知性は、その瞬間に機能を停止する。

ゆえに、社会がゲティア問題を解決しないのは、怠慢や不徹底によるのではない。それは構造的必然である。社会的停止構造が機能するためには、ゲティア問題を解決しないことが必要なのである。完全な正当化の放棄は、社会的知性の動作条件である。

\bigskip


\subsection{知識の再定義 — 行動可能化のための暫定ブロック}


以上から、知識の動力学的再定義が導かれる。

伝統的定義「知識 = 正当化された真なる信念」は、知識を $L_F$ 内部の性質(正当化)と外部実在との対応(真理)の合成として定義する。本論文の枠組みでは、知識はこのようには定義されない。

\textbf{定義 10.1 (動力学的知識)}

\begin{quote}
知識とは、行動可能化のために $\Phi_{\text{net}}$ において暫定的に採用された停止ブロックである。
\end{quote}


すなわち知識は、真理との対応によってではなく、行動を可能にする暫定的停止としての機能によって定義される。知識は $L_E$ が供給する停止であり、$\Phi_{\text{net}}$ において暫定採用され、改訂可能に保たれる。

この定義において、ゲティア事例は問題ではなくなる。ゲティア事例は「正当化され真だが知識でない」事例として逆説的に見えたが、動力学的知識の定義では、知識は正当化と真理の合致によって定義されないため、両者が合致してもしなくても、知識の成立は $L_E$ の停止が安定し行動可能化に資するかどうかによる。ゲティア事例の逆説は、知識を正当化+真理として定義したことの産物であって、動力学的知識には現れない。

\bigskip


\subsection{真理に依存しない操作的安定性}


動力学的知識の核心は、真理に依存しない操作的安定性である。

知識が行動可能化のための暫定ブロックであるなら、その安定性は真理との対応にではなく、操作的機能に依存する。ある停止ブロックが知識として安定するのは、それが真であるからではなく、それが行動を可能にし、$\Phi_{\text{net}}$ において改訂可能に維持されるからである。

ここで \S{}1.1 で予告した区別が回収される。閉包度 $C$ は安定性測度であって真理測度ではない。$C$ が高いことは、停止が動力学的に安定であることを意味するのであって、その停止が真理に対応することを意味しない。$L_E$ が供給する停止の $C$ は、操作的安定性を測るのであって、正しさを測らない。

厳密な正しさの要求は、操作的安定性を破壊する。完全な正しさを停止の条件とすれば、\S{}10.3 で見た通り、停止は不可能になり、推論は無限後退に戻る。操作的安定性は、正しさの要求を緩和することによってのみ成立する。これは妥協ではなく、停止構造が機能するための構造的条件である。

\subsubsection{プラグマティズムとの区別}


知識を行動可能化との関係で捉える本論文の立場は、プラグマティック認識論(James, Dewey らの系譜)と帰結において響き合う。ただし本論文はこれを既存思想から借りるのではなく、VED の動力学的経路——非閉包定理(\S{}3)と創発停止層(\S{}5)——から独立に導出している。

より重要なのは、本論文の立場が俗流的に理解されたプラグマティズムとは異なることである。俗流理解では、プラグマティズムはしばしば「役に立てば真である」「真理とは有用性にほかならない」という、真理の有用性への還元として戯画化される。本論文はこの主張をしない。

本論文は真理を否定も放棄もしていない。本論文が述べるのは、停止($L_E$)の成立条件が真理との対応ではなく操作的安定性にある、という構造的位置づけのみである。真理は依然として存在しうる。しかしそれは、推論層 $L_F$ の正当化とも、創発停止層 $L_E$ の停止とも、別の事柄である。「真理 = 有用性」と述べるのではなく、「停止の成立は真理判定とは別の動力学的事柄である」と述べているのである。$C$ が真理測度でないことは、真理が存在しないことを意味するのではなく、$C$ が測っているのが真理ではないことを意味するに過ぎない。

この区別は重要である。本論文は真理概念を解体するのではなく、停止の動力学を真理判定から分離する。両者を混同すると、本論文は「真理を有用性に還元する相対主義」として誤読されうるが、それは本論文の主張ではない。

\bigskip


\subsection{誤り許容性の必然と $\mu_\theta$ との同値}


操作的安定性が真理に依存しないことから、誤り許容性が必然的に導かれる。

誤り許容性とは、誤った停止を許容する構造的余地である。$L_E$ が供給する停止は真理を保証しないため、誤った停止が生じうる。社会的停止構造は、この誤りを許容する。誤った停止も暫定ブロックとして採用され、後に $\Phi_{\text{net}}$ の相互干渉によって改訂される(\S{}9)。

誤りを許容しない停止構造は機能しない。誤りを一切許さないためには完全な正当化が必要であり、それは到達不可能だからである(\S{}3)。誤りなき社会的停止構造は、判断生成それ自体を停止する。誤り許容性は、停止構造が判断を生成し続けるための必要条件である。

ここで、本論文の第二論理線との重要な接続が現れる。誤り許容性は、閾値移動能力 $\mu_\theta$ の存在と構造的に同値である。

論証は次の通りである。$\mu_\theta$ は運動可能領域 $\Omega_\theta$ を組み替える能力であった(\S{}2.7)。領域が組み替え可能であるということは、現在の停止が確定的・最終的でないことを意味する。確定的でない停止は、誤りうる停止である。逆に、誤りを許容する停止構造は、停止が改訂されうることを前提とし、それは $\theta$ が動きうること——$\mu_\theta$ の存在——を要求する。

したがって、ゲティア構造の排除不能性(誤りの構造的可能性)は、$\mu_\theta$ の存在の必然的代償である。$\mu_\theta$ を持つ——閾値を動かし、運動可能領域を組み替えられる——ことと、$C$ と真理を同一視できない——誤りが構造的に排除できない——ことは、同じ構造の二つの側面である。

これは深い帰結を持つ。サピエンス的知性が $\mu_\theta$ を内発的に持つこと(\S{}5.4)と、サピエンス的知性がゲティア問題を抱えること——正しさによっては知識を閉じられないこと——は、切り離せない。閾値を動かせる知性は、必然的に誤りうる知性である。完全に正しい知性は、閾値を動かせない硬直した知性であり、それは \S{}11 の硬直域に対応する。

\bigskip


\subsection{ゲティア問題の最終的役割}


ゲティア問題の役割を総括する。

ゲティア問題は、認識論における未解決の難問ではない。それは、知性が正しさのみでは閉じないことの証明である。

\S{}3 が「$L_F$ は構造的に閉じない」を示し、\S{}5 が「$L_E$ が外部依存的に停止を供給する」を示した。ゲティア問題は、この構造が認識論において必然的に生む帰結である。知識を正当化+真理として閉じようとする試みがことごとく失敗するのは、知識の成立が $L_F$ 内部にも、真理との対応にもなく、$L_E$ の操作的安定性にあるからである。

ゲティア問題は、こうして、本論文の第一中心(\S{}3)と第二中心(\S{}5)を認識論において確証する。知性は正しさによって閉じるのではない。知性は $L_E$ の改訂可能な停止によって、外部依存的に、誤りを許容しながら、操作的に安定する。ゲティア問題は、この構造の認識論的な指紋である。

\bigskip


\subsection{次章への橋渡し}


本章で、$C$ が安定性測度であり真理測度でないこと、誤り許容性が $\mu_\theta$ と同値であることを示した。ここから、知性の安定性に関する定量的な問いが開かれる。

$\mu_\theta$ を持つ知性は誤りうる。では、$\mu_\theta$ の作用——閾値移動、運動可能領域の組み替え——が、速度・解像度・停止硬度といったパラメータとどう関係するとき、知性は安定し、どう関係するとき不安定化するのか。誤り許容性が失われ、$\mu_\theta$ が機能しなくなるのはどのような条件下か。

次章 \S{}11 で扱う不安定領域は、この問いに答える。速度 $v_{\text{inf}}$・解像度 $r_{\text{res}}$・停止硬度 $\eta$ の三軸相図において、知性の安定運用領域と不安定領域(発散域・硬直域)を区別し、$\mu_\theta$ と $L_E$ がどの条件下で機能し、どの条件下で崩壊するかを分析する。

\bigskip


\subsection{要約}


本章で示したのは次の点である:

\begin{enumerate}
\item ゲティア問題は周辺的反例ではなく、知識が正当化と真理の合成では定義できないことの構造的暴露である
\item ゲティア問題は原理的に解決不可能である。$L_F$ 内部に停止条件がないため(\S{}3)、第四条件の追加は無限後退する
\item 社会はゲティア問題を解決しない。完全な正当化を要求すれば社会的知性は即座に停止するため、解決しないことが構造的必然である
\item 知識は「行動可能化のために $\Phi_{\text{net}}$ で暫定採用された停止ブロック」と再定義される(定義 10.1)
\item 知識の安定性は真理でなく操作的安定性による。$C$ は安定性測度であって真理測度ではない(\S{}1.1 の回収)。ただしこれは真理の有用性への還元(俗流プラグマティズム)ではなく、停止の動力学を真理判定から分離するものである
\item 誤り許容性は必然であり、閾値移動能力 $\mu_\theta$ の存在と構造的に同値である。閾値を動かせる知性は必然的に誤りうる
\item ゲティア問題は、知性が正しさのみでは閉じないことの証明であり、第一中心と第二中心の認識論的確証である
\end{enumerate}


次章 \S{}11 では、知性の安定性を速度・解像度・停止硬度の三軸相図において分析する。

\bigskip
